\section{Discussion and conclusions} \label{Con}

We have applied the local variance estimator to the latest Planck release data set (PR4) in both temperature and polarization, performing a frequentist analysis of the isotropy assumption. In temperature, we have not found any simulation from a set of 600 with an amplitude as large as the one observed in the data for $4^{\circ}$, $6^{\circ}$, and $8^{\circ}$ discs. This means that the data exhibits a dipolar-like behaviour across the sky at a level below 0.17\% for that range of angular scales. We have performed the analysis at two different resolutions, $N_{side} = 64$ and $N_{side} = 2048$. For low resolution maps the direction points towards (l, b) = ($208^{\circ}$, $-15^{\circ}$), while for full resolution the direction is (l, b) = ($205^{\circ}$, $-20^{\circ}$). These results are in good agreement with previous reports, and can not be explained by residual systematics since a similar detection is found in WMAP.

Regarding polarization, we have implemented an alternative inpainting approach in order to minimize some effects that affect the reconstruction of the E-mode map from the measured Stokes parameters. This technique works optimally under certain conditions as will be shown in a future work (Gimeno-Amo et al. in preparation) that are not completely satisfied in the Planck data. The main problem is the low number of realistic noise simulations, which does not allow one to compute in an accurate way the noise covariance matrix. However, we have checked in this work that the inpainting is not biasing the parameters of the analysis and that also improves significantly the results compared with a simple masking approach.

First of all, we have applied the inpainting technique to the PR3 data set getting results that are in good agreement with the ones obtained by the Planck collaboration using a different inpainting algorithm. 
The small differences found can be explained by the large intrinsic dispersion of the estimated dipole direction, as well as differences in the considered mask and inpainting algorithm.
In particular we tested that both the p-values and the directions are affected by the considered mask, with the exception of Sevem,
whose p-value is always below 1\%.

For PR4, we have performed the analysis only for Sevem. We get a p-value that lies between 1\%-3\% for sky fractions in the range from 67\% to 59.4\%. For smaller $f_{\rm sky}$ the p-value increases rapidly, getting up to 7.4\% for $f_{\rm sky}$ = 57\%. If the assymetry is of cosmological origin, this behaviour is somehow expected due to the loss of information when the number of available pixels is reduced. 
However, we can not discard that this effect is related to the reduction of foreground contamination in the CMB map when increasing the mask, that one would expect if residuals were the cause of the observed assymetry. Finally, the results are probably even more sensitive to the sky fraction and geometry of the mask, due to the fact that the estimated direction of the dipole is close to its boundary.


It is also interesting to note that the amplitude of the dipole estimated for intensity ($\sim$7\%) and polarization ($\sim$9\%) are of similar level and point in a similar direction. If the intensity dipole were simply a statistical fluke such that the standard model still holds, we could expect a similar dipolar pattern
in the E-map, 
taking into account the correlation between T and E. 
However, since this correlation is small (typically around 10 per cent), a much lower amplitude in the dipole in polarization would be expected and, therefore, this seems inconsistent with the values found. Even more, we have checked with E-mode maps constrained to intensity modulated simulations, that this correlation is expected to increase the p-value of the alignment by only about $0.5\%$ when realistic E2E noise is included. In any case, given the large uncertainties associated to the estimated amplitudes, improved polarization data is needed before any further conclusion can be established.

We may wonder whether the difference found between the p-values of PR3 and PR4 (0.22\% for PR3 versus 2.8\% for PR4 for Sevem for our reference case) is significant. There are several effects that can cause this difference. First of all, especially since we are in the tail of the distribution, the limited number of simulations may play a role. Indeed, if we consider the 95 per cent range for the p-values of both pipelines, they overlap, even if this range only takes into account the dispersion due to the different possible splits. Additional uncertainty would come from the estimator itself, including the inpainting. Another point to take into account is that we are using a different mask. Although we fix the same threshold in the maximum error of the reconstruction of the E-mode, this leads to a different sky fraction for PR3 and PR4, which can also influence the found p-value. Nevertheless, there are other reasons related to the data themselves that could lead to these results. In particular, if the anomaly were due to systematic effects, it would be expected to find a lower significance for PR4, due to the improvement in this regard in the data. However, if the anomaly were of cosmological origin, the significance could also be reduced due to the presence of the transfer function that affects the low multipoles in polarization for PR4. Therefore, not clear conclusion can be derived from these results, except that a modest detection of asymmetry is still present in the E-mode polarization maps.

Unfortunately, the sensitivity of the data is still not sufficiently good to obtain robust results. In particular, from the model-dependent test done in this paper we see how the model parameters present large uncertainties. Furthermore, with the Planck polarization data set would be only possible to claim a detection with a p-value lower than 1\% in 95\% of the cases if the amplitude of the modulation were at the level of 16\%, well above the values found. Therefore, significant improvement in the sensitivity and systematics of the polarization data is needed on large angular scales before a robust result can be achieved in relation to possible violation of isotropy. Future experiments, such as LiteBIRD, that will be able to measure the cosmological E-modes at the cosmic-variance limit over 
a very large fraction of the sky, are expected to provide key information on the HPA.
